\documentclass[11pt,a4paper]{article}
\usepackage[utf8]{inputenc}
\usepackage[T1]{fontenc}
\usepackage[spanish]{babel}
\usepackage{amsmath,amssymb,amsthm}
\usepackage{geometry}
\geometry{margin=2.5cm}

\newtheorem{definition}{Definici\'on}
\newtheorem{proposition}{Proposici\'on}
\newtheorem{remark}{Observaci\'on}

\title{Auditor\'ia Estructural OVRP de la Formalizaci\'on en Lean 4 de la Hip\'otesis de Riemann y Escaneo Espectral en la L\'inea Cr\'itica}
\author{Mariano Panzano Caball\'e \\ \small Oasis Sovereign Architecture Research, Barcelona}
\date{Noviembre de 2026}

\begin{document}
\maketitle

\begin{abstract}
Aplicamos el \emph{Oasis Verifiable Research Pipeline} (OVRP, DOI: \texttt{10.5281/zenodo.24443531}) para auditar la estructura matem\'atica y la huella axiom\'atica de la formalizaci\'on en Lean~4 de la Hip\'otesis de Riemann publicada en noviembre de 2026 (175.187 m\'odulos, huella est\'andar \texttt{[propext, Classical.choice, Quot.sound]}), contrast\'andola con el modelo f\'isico-computacional de consenso distribuido y correspondencia primo-hologr\'afica de Oasis Swarm. Asimismo, documentamos un barrido espectral determinista sobre $\sigma \in \{0.3, 0.4, 0.5, 0.6, 0.7\}$ en $t \in [13.5, 38.5]$ y formalizamos las invariantes de simetr\'ia funcional alrededor de $\operatorname{Re}(s) = 1/2$ en Lean~4.36.0 (\texttt{0 sorry}).
\end{abstract}

\section{Resultados de la Auditor\'ia Estructural OVRP}
Mediante extracci\'on determinista en memoria RAM sobre la reconstrucci\'on anal\'itica de 137 p\'aginas (\texttt{DOI: 10.5281/zenodo.24433694}) y los registros de verificaci\'on de compilaci\'on (\texttt{Main.olean}, SHA-256 \texttt{9146dfb4\dots}), se constata:
\begin{enumerate}
    \item La demostraci\'on formal utiliza directamente la definici\'on can\'onica \texttt{RiemannHypothesis} de Mathlib y no introduce axiomas personalizados.
    \item Su arquitectura matem\'atica se apoya en series theta (205 menciones), levantamientos theta sobre grupos unitarios totalmente definidos, cohomolog\'ia de Galois y teor\'ia de Iwasawa, sin recurrir a operadores de flujo de calor ni cotas entr\'opicas expl\'icitas.
\end{enumerate}

\section{Evidencia Emp\'irica en la L\'inea Cr\'itica y Frontera Epistemol\'ogica}
En el modelo de consenso de red de primos $\Psi(s) = \prod_p (1 - p^{-s})^{-1}$, el escaneo num\'erico (\texttt{tests/riemann\_spectral\_scan.py}) confirma m\'inimos profundos de resonancia exclusivamente sobre $\operatorname{Re}(s) = 0.5$ en $t \in [13.5, 38.5]$, mientras que las l\'ineas paralelas $\sigma \in \{0.3, 0.4, 0.6, 0.7\}$ carecen de m\'inimos estables.

\begin{remark}[Frontera Expl\'icita]
El modelo primo-hologr\'afico y el barrido num\'erico en intervalos finitos proporcionan soporte emp\'irico y computacional coherente dentro de Oasis Swarm, pero no constituyen por s\'i mismos una demostraci\'on anal\'itica incondicional ni sustituyen a una prueba formal completa en teor\'ia anal\'itica de n\'umeros.
\end{remark}

\end{document}
